% Thai edition. Build with XeLaTeX (Thai line breaking needs ICU): xelatex paper_th.tex
\documentclass[11pt]{article}
\usepackage[margin=0.8in]{geometry}
\usepackage{fontspec}
\setmainfont{Sarabun}[Path=fonts/, Extension=.ttf, UprightFont=*-Regular, BoldFont=*-Bold,
  ItalicFont=*-Italic, BoldItalicFont=*-BoldItalic]
\setmonofont{Menlo}[Scale=0.85]
\XeTeXlinebreaklocale "th"
\XeTeXlinebreakskip = 0pt plus 1pt
\linespread{1.18}
\usepackage{booktabs}
\usepackage{graphicx}
\usepackage{hyperref}
\usepackage{xcolor}
\usepackage{float}
\usepackage{caption}
\usepackage{array}
\hypersetup{colorlinks=true,linkcolor=blue!55!black,citecolor=blue!55!black,urlcolor=blue!55!black}
\renewcommand{\abstractname}{บทคัดย่อ}
\renewcommand{\refname}{เอกสารอ้างอิง}
\renewcommand{\figurename}{รูปที่}
\renewcommand{\tablename}{ตารางที่}
\input{generated/macros.tex}
\title{Raise, Limp หรือ Fold ที่ 20 Big Blinds:\\
ทดสอบการจัดกลุ่มมือเริ่มต้น PLO ของ Jeff Hwang\\และเสนอกลุ่มมือใหม่ที่ตรงกับ Solver}
\author{Nuttakit Kundum}
\date{กันยายน 2026}

\begin{document}
\maketitle

\begin{abstract}
การแบ่งมือเริ่มต้นของ Jeff Hwang เป็น Premium, Speculative, Marginal และ Trash เป็นวิธีเลือกมือที่นิยมใน
pot-limit Omaha งานนี้ถามว่าการแบ่งแบบนั้นอธิบายสิ่งที่ solver ทำกับการตัดสินใจแรกได้ดีแค่ไหน ในเกม 6 คนที่ stack
20 big blinds เมื่อยังไม่มีใครเข้า pot และผู้เล่นเลือกได้ว่าจะ fold, limp หรือ raise ขนาด pot เราแก้เกมที่ bet
ขนาด pot ได้ทุก street และ re-raise ได้จนถึง all-in โดยใช้การจัดกลุ่มมือที่ไม่เอามือต่าง tier ของ Hwang มาไว้กลุ่ม
เดียวกัน แล้วตรวจมือจริงทีละมือโดยเทียบมูลค่าของการ fold, limp และ raise ผลมีสามข้อ ข้อแรก tier ของ Hwang
เรียงความแข็งของมือได้ถูก แต่ใช้เป็นกฎที่ไม่ขึ้นกับตำแหน่งไม่ได้ มือ Premium \PremFoldUTG\% ของคอมโบทั้งหมด
ควร fold ที่ UTG อย่างมีนัยสำคัญ และทั้งหมดเป็นมือสาย straight ไพ่ต่ำที่ไม่ใช่ double-suited ข้อสอง
การตัดสินใจแรกขึ้นกับจำนวนไพ่ตั้งแต่ T ขึ้นไปมากที่สุด ตามด้วยการมี A ที่มีดอกคู่ และระดับของคู่ ส่วนความต่อเนื่อง
ของไพ่สาย straight ซึ่งเป็นแกนของระบบ Hwang แทบไม่มีผลด้วยตัวเองที่ความลึกนี้ ข้อสาม archetype ใหม่
\ArchCount\ กลุ่มที่สร้างจากปัจจัยเหล่านี้อธิบายอัตรา fold ที่ UTG ของ solver ได้ $R^2 = \RsqArchUTG$ เทียบกับ
$R^2 = \RsqTierUTG$ ของ tier ของ Hwang และลด EV ที่เสียจากการตัดสินใจแรกผิดจาก \LossHwangUTG\ เหลือ
\LossArchUTG\ big blinds ต่อครั้ง บนมือที่แยกไว้ทดสอบ
\end{abstract}

\section{คำถาม}
Hwang แบ่งมือสี่ใบที่เล่นได้ตามโครงสร้าง ได้แก่ ไพ่ใหญ่และ Broadway wrap ที่นำด้วย A, มือสาย straight,
มือ A มีดอก, มือคู่บวกไพ่อื่น, มือ AA และมือ marginal \cite{hwang} แล้วให้ระดับแต่ละรูปแบบเป็น Premium,
Speculative, Marginal หรือ Trash ตัวจัดกลุ่มของโปรเจกต์นี้ทำตามระดับนั้นจากอันดับไพ่ ช่องว่าง และดอก โดยจัดมือ
5.49\% จาก 270{,}725 มือเป็น Premium, 13.83\% เป็น Speculative, 19.91\% เป็น Marginal และ 60.76\% เป็น Trash

จุดเริ่มต้นคือข้ออ้างแคบ ๆ ว่ามือ Premium ไม่ควร fold เลยเมื่อยังไม่มีใครเข้า pot และเราเป็นคนแรกที่ตัดสินใจ
เราขยายเป็นสามคำถาม: solver raise, limp หรือ fold มือไหนในแต่ละตำแหน่ง, การให้ระดับของ Hwang ขัดกับ solver
ตรงไหน และ archetype ชุดเล็ก ๆ ชุดใหม่อธิบาย solver ได้ดีกว่าหรือไม่

คำแนะนำของ Hwang เองมีเงื่อนไข เขาเขียนว่ามือระดับ Premium ``มักเล่นได้จากทุกตำแหน่งและในสถานการณ์ส่วนใหญ่''
ว่าควร ``limp จากตำแหน่งหน้าเป็นปกติ'' และว่ามือ Marginal เป็น ``มือตำแหน่งหลัง ลงขั้นต่ำเท่านั้น'' \cite{hwang}
หนังสือของเขามุ่งไปที่เกม cash ที่ stack ลึกและเล่นกันหลายคน stack 20 big blinds เป็นอีกสถานการณ์หนึ่ง
และส่วนหนึ่งของงานนี้คือดูว่าความคิดไหนของเขายังใช้ได้

\section{เกมและ solver}
ผู้เล่น 6 คน blinds 0.5/1 stack คนละ 20 big blinds ไม่มี ante ไม่มี rake และวัดผลเป็น chip EV ทุกการ raise
เป็นขนาด pot และไม่เกิน all-in ก่อน flop raise ได้สามครั้ง (open, three-bet และ four-bet ซึ่งแทบเท่ากับ shove)
ใน flop, turn และ river bet ขนาด pot ได้หนึ่งครั้งและ raise ได้อีกสองครั้ง side pot, pot แบ่งเท่า และกฎ showdown
ของ Omaha ที่ต้องใช้ไพ่ในมือสองใบกับไพ่กลางสามใบคำนวณแบบตรงทุกกรณี tree สาธารณะมี 474{,}658 node และ
232{,}258 จุดตัดสินใจ

\paragraph{การจัดกลุ่มมือ} ก่อน flop กลุ่มโครงสร้าง 572 กลุ่มถูกแยกตาม tier ของ Hwang เป็น 780 กลุ่ม ทำให้
ไม่มีกลุ่มไหนปน tier กัน ใน flop และ turn มือถูกจัดกลุ่มตามสัดส่วนของไพ่สองใบบนกระดานที่ชนะมือเรา (10 ช่วง
โดยแยกช่อง nut ไว้ต่างหาก) ตาม flush draw (ไม่มี, ไม่ใช่ nut, nut) และตามจำนวนอันดับไพ่ใบถัดไปที่ทำ straight ได้
(0, 1, 2, 3 ขึ้นไป) รวม 120 กลุ่มต่อ street ส่วน river ใช้ความแข็งของมืออย่างเดียว

\paragraph{การฝึก} ใช้ external-sampling Monte Carlo CFR พร้อม linear discounting อัปเดตแบบขนานบนคอมพิวเตอร์เครื่องเดียว
ฝึกสอง seed ที่แยกกันอิสระ seed ละ 100 นาที ได้ \Deals\ ล้านดีลต่อ seed ดีลละหก traversal
ความถี่การตัดสินใจแรกระดับ tier ของสอง seed ต่างกันไม่เกิน 3.1 จุดในทุกตำแหน่ง ยกเว้น small blind

\paragraph{การตรวจสอบ} ตัวประเมินมือ Omaha ให้ผลตรงกับ \texttt{phevaluator} ทุกกลุ่มถูกตรวจว่ามี tier เดียว
ชิปรวมคงที่ในทุก terminal และกรณีที่คำนวณด้วยมือได้ (ทุกคน fold ให้คน open ได้ +1.5 big blinds) ตรงเป๊ะ
เมื่อสุ่มมือ 150 มือต่อตำแหน่ง การเปลี่ยนการตัดสินใจแรกให้ดีขึ้นได้กำไรเฉลี่ยไม่เกิน 0.05 big blinds
แปลว่าการตัดสินใจแรกเกือบสอดคล้องกันเองแล้ว ตัวเลขนี้เป็นขอบล่างของ exploitability ไม่ใช่ใบรับรอง

\section{หลักฐานสองแบบ}
\paragraph{ความถี่จาก solver} สำหรับมือทั้ง 16{,}432 คลาสที่แยกตามความสมมาตรของดอก เราบันทึกความน่าจะเป็น
fold, limp และ raise ที่ตำแหน่งเข้าเล่นคนแรกแต่ละตำแหน่ง เฉลี่ยจากสอง seed ข้อมูลนี้ครอบคลุมทุกมือ แต่ solver
เลือกทีละกลุ่ม มือสองมือในกลุ่มเดียวกันจึงได้สัดส่วนเดียวกันเสมอ

\paragraph{การตรวจทีละมือ} สำหรับมือจริงหนึ่งมือที่จุดเข้าเล่นคนแรก เราประมาณมูลค่าของการ fold, limp และ raise
เมื่อทุกคนรวมถึงตัวเราเล่นตามกลยุทธ์ที่แก้ได้หลังจากนั้น แต่ละดีลถ่วงน้ำหนักด้วยความน่าจะเป็นที่ทุกตำแหน่งก่อนหน้า
fold มือหนึ่ง \emph{ควรเข้าเล่นอย่างชัดเจน} ถ้า limp หรือ raise ดีกว่า fold อย่างน้อยสามเท่าของ standard error,
\emph{ควร fold อย่างชัดเจน} ถ้าทั้งสองแย่กว่า fold อย่างน้อยสามเท่าของ standard error และ \emph{ก้ำกึ่ง}
ในกรณีอื่น เราตรวจคลาส Premium ครบทั้ง 922 คลาส และสุ่มมือ 400 มือจากแต่ละ tier ที่เหลือ ในทั้งห้าตำแหน่ง
รวม \AuditRows\ แถวจาก \AuditClasses\ คลาส ถ่วงน้ำหนักให้แทนประชากรทั้งหมด เมื่อเทียบสอง seed มีมือ
Premium ที่ UTG \SeedContra\ จาก \SeedPairs\ มือที่ได้ผลชัดเจนสวนทางกัน การตรวจทีละมือแยกมือที่อยู่กลุ่มเดียวกัน
ได้ เช่น KK32 single-suited ได้ +0.30 big blinds เมื่อ limp จาก UTG แต่ TT32 เสีย 0.93 ทั้งที่ solver
เล่นสองมือนี้เหมือนกัน

\section{แต่ละ tier fold บ่อยแค่ไหน}
\begin{table}[H]
\centering
\caption{สัดส่วนการตัดสินใจแรกจาก solver เป็น fold/limp/raise (\%) เฉลี่ยสอง seed}
\small
\input{generated/table_tiers_th.tex}
\end{table}

\begin{table}[H]
\centering
\caption{ผลตรวจทีละมือ: สัดส่วนคอมโบที่ควรเข้าเล่นชัดเจน / ก้ำกึ่ง / ควร fold ชัดเจน (\%)}
\small
\input{generated/table_verdicts_th.tex}
\end{table}

\begin{figure}[H]
\centering
\includegraphics[width=\linewidth]{generated/fig_verdicts_th.pdf}
\caption{ผลตรวจทีละมือแยกตาม tier และตำแหน่ง}
\end{figure}

tier ยังเรียงลำดับเดิมในทุกตำแหน่ง Premium ถูก fold น้อยที่สุด Trash มากที่สุด และทุก tier ถูกเล่นมากขึ้นเมื่อ
ตำแหน่งขยับเข้าใกล้ button ที่ UTG มือ Premium \PremEnterUTG\% ของคอมโบควรเข้าเล่นและ \PremFoldUTG\%
ควร fold ส่วนที่ button สัดส่วนที่ควร fold ลดเหลือ \PremFoldBTN\% ขณะที่ Trash ควร fold ชัดเจน
\TrashFoldUTG\% ของคอมโบที่ UTG ตำแหน่ง small blind ต่างจากที่อื่นในทุก tier เพราะเจอแค่ big blind และจ่าย
ไปแล้วครึ่ง big blind solver จึง limp มือจำนวนมาก รวมถึง Trash 45\% และแทบไม่มีมือไหนที่ควร fold ชัดเจน
มือ Marginal ส่วนใหญ่และมือ Speculative จำนวนมากอยู่ในกลุ่มก้ำกึ่ง แปลว่าที่ 20 big blinds มือพวกนี้ใกล้เสมอตัว
มากกว่าเป็นมือที่ต้อง fold ชัด ๆ

\section{อะไรตัดสินว่าจะ raise, limp หรือ fold}
ตารางถัดไปแบ่งมือทั้งหมดตามลักษณะทีละอย่าง ``เล่น'' คือความน่าจะเป็นที่ solver limp หรือ raise ส่วน ``EV
เข้าเล่น'' คือมูลค่าของทางที่ดีกว่าระหว่าง limp กับ raise ลบด้วยการ fold หน่วยเป็น big blinds จากการตรวจทีละมือ

\begin{table}[H]
\centering
\caption{อัตราการเล่นและ EV เข้าเล่นแยกตามลักษณะทีละอย่าง}
\small
\input{generated/table_drivers_th.tex}
\end{table}

\paragraph{ไพ่ใหญ่มาก่อน} มือที่มีไพ่ตั้งแต่ T ขึ้นไปศูนย์หรือหนึ่งใบแทบไม่ถูกเล่นจากตำแหน่งหน้า และเสียเกือบหนึ่ง
big blind เมื่อเล่น มือที่มีไพ่ใหญ่สามใบถูกเล่น 69\% ที่ UTG และ 96\% ที่ button ส่วนสี่ใบแทบเล่นทุกครั้ง
ไม่มีลักษณะเดี่ยวอื่นที่แยกมือได้ชัดเท่านี้

\paragraph{ดอกของ A สำคัญกว่าดอกที่สอง} การมี A ที่มีดอกคู่ยก EV เข้าเล่นที่ UTG จาก -0.45 (ดอกที่ดีที่สุด
นำด้วย T ถึง K) เป็น +0.24 ดอกที่สองช่วยได้ (double-suited +0.12 เทียบ single-suited -0.53)
ส่วนมือที่ดอกดีที่สุดนำด้วย 9 หรือต่ำกว่าเล่นแทบไม่ต่างจากมือ rainbow

\paragraph{คู่ขึ้นกับระดับ} AA ถูกเล่นเสมอและมีค่าราวสาม big blinds KK และ QQ ได้กำไรจากทุกตำแหน่ง
JJ และ TT ใกล้เสมอตัวที่ UTG (-0.12) และได้กำไรเฉพาะตำแหน่งหลัง คู่เล็กควร fold ยกเว้นที่ small blind

\paragraph{ความต่อเนื่องของไพ่อย่างเดียวช่วยน้อย} มือที่ไพ่ทั้งสี่อยู่ในช่วงห้าอันดับถูกเล่นที่ UTG แค่ 35\% และ
EV เข้าเล่นยังติดลบ rundown มีค่าเมื่อมีไพ่ใหญ่และดอกมาด้วย ซึ่งเป็นปฏิสัมพันธ์ที่ archetype ของเราใช้

\paragraph{Limp หรือ raise} AA และมือที่มีไพ่ใหญ่ตั้งแต่สามใบพร้อม A มีดอกหรือ double-suited ถูก raise
บ่อยกว่า limp มากในทุกตำแหน่ง และ raise ดีกว่า limp ราว 0.15--0.3 big blinds คู่ใหญ่ที่ไม่มี A มีดอกถูก limp
ราว 60\% ที่ UTG แม้มูลค่าของ limp กับ raise ต่างกันแค่ราว 0.1 big blinds rundown ไพ่ใหญ่เป็นกลุ่มเดียวที่ limp
จาก UTG ดีกว่าชัดเจน (+0.22 เทียบ +0.12) ซึ่งตรงกับคำแนะนำของ Hwang ที่ให้มือ draw เล่นหลายคนจากตำแหน่งหน้า
ส่วนที่ small blind มือที่เข้าเล่นถูก raise บ่อยกว่า limp ยกเว้นกลุ่มที่อ่อนที่สุด

\section{Hwang ขัดกับ solver ตรงไหน}
ตารางต่อไปนี้แสดงกลุ่มใหญ่ที่สุด (tier, รูปแบบ และดอกตามแบบของ Hwang) ที่คำแนะนำของ Hwang กับ solver ชี้ไปคนละ
ทางที่ UTG เราตีความคำแนะนำของ Hwang ว่า raise มือ Premium, limp มือ Speculative, เล่นมือ Marginal เฉพาะ
cutoff, button และ small blind และ fold มือ Trash

\begin{table}[H]
\centering
\caption{ความขัดแย้งที่ใหญ่ที่สุดที่ UTG EV เข้าเล่นเป็นค่าเฉลี่ยจากการตรวจทีละมือเมื่อมีข้อมูล}
\small
\resizebox{\linewidth}{!}{\input{generated/table_mismatch_th.tex}}
\end{table}

ความขัดแย้งส่วนใหญ่มาจากสามรูปแบบ
\begin{itemize}
\item \textbf{มือสาย straight และมือคู่บวกไพ่ต่อถูกให้ค่าสูงเกินไป} มือคู่บวกไพ่ต่อแบบ single-suited (เช่น 9-9-8-7)
และมือ straight มีช่องว่างที่ไม่มีไพ่ใหญ่ เป็น Speculative หรือ Premium ในระบบของ Hwang แต่ที่ UTG เสียราว
0.35--0.7 big blinds เทียบกับ fold มือพวกนี้ต้องพึ่ง stack ลึกและ implied odds
\item \textbf{มือไพ่ใหญ่ที่มีไพ่ห้อยถูกให้ค่าต่ำเกินไป} คู่ใหญ่ที่มีไพ่ห้อย Broadway สามใบที่มีไพ่ห้อย และมือ A มีดอก
อ่อน ๆ แบบ double-suited เป็น Marginal สำหรับ Hwang แต่ solver เล่นมือเหล่านี้ส่วนใหญ่จาก UTG และแบบ
double-suited ได้กำไร (+0.17 ถึง +0.66) ที่ 20 big blinds ความแข็งของไพ่สูงมักได้ผลตอน showdown หรือด้วยการ
ทุ่มครั้งเดียวบน flop ไพ่ห้อยจึงเสียน้อยกว่าตอน stack ลึก
\item \textbf{การนับดอกหยาบเกินไป} ตัวจัดกลุ่มถือว่าไพ่สองใบดอกเดียวกันคือ ``มีดอก'' แต่ solver แยก A มีดอก,
K ถึง T มีดอก และดอกต่ำออกจากกันชัดเจน และเล่นมือ three-flush น้อยเกือบเท่ามือ rainbow (9\% และ 6\% ที่ UTG
เทียบกับ 20\% ของมือ single-suited)
\end{itemize}

\begin{figure}[H]
\centering
\includegraphics[width=0.52\linewidth]{generated/fig_rundowns_th.pdf}
\caption{มือสาย straight ระดับ Premium ที่ UTG: EV เข้าเล่นที่ดีที่สุดลบ fold หน่วย big blinds แยกตามอันดับไพ่และ
แบบดอก (ds = double-suited, ss = single-suited, 3f = three-flush, mono = monotone)}
\end{figure}

รูปนี้แสดงกรณี Premium อย่างละเอียด rundown แบบ double-suited ใกล้เสมอตัวหรือดีกว่า และ rundown ที่นำด้วยไพ่ใหญ่
สองใบ (J-T-9-8, Q-J-T-9, K-Q-J-9) เล่นได้เมื่อมีดอก แต่ rundown แบบ single-suited, three-flush และ monotone
ตั้งแต่ 6-5-3-2 ถึง 9-8-7-6 เสีย 0.4--1.3 big blinds เทียบกับ fold มือเหล่านี้ได้ป้าย Premium เพราะตัวจัดกลุ่มให้ระดับ
มือสาย straight จากช่องว่าง และถือว่ามีดอกอะไรก็พอ ขณะที่ที่ความลึกนี้ equity ของมันเมื่อถูก call ต่ำเกินไป

\section{Archetype ใหม่ที่ตรงกับ solver}
ปัจจัยข้างต้นชี้ว่าควรจัดกลุ่มมือตามจำนวนไพ่ใหญ่ ดอกของ A ดอกที่สอง และระดับของคู่ โดยใช้ความต่อเนื่องของไพ่
เฉพาะเมื่อมีไพ่ใหญ่ด้วย ตารางแรกด้านล่างเป็นกฎที่ตรวจตามลำดับ ตารางที่สองเป็นการตัดสินใจแรกที่แนะนำในแต่ละ
ตำแหน่ง ซึ่งมาจากค่าเฉลี่ยของ limp และ raise ที่ตรวจทีละมือ ช่องที่เขียนว่า ``ก้ำกึ่ง'' คือทางเข้าเล่นที่ดีกว่าอยู่ห่าง
จาก fold ไม่ถึง 0.10 big blinds

\begin{table}[H]
\centering
\caption{กฎของ archetype ไพ่ใหญ่คือ T ขึ้นไป คู่ใหญ่นับเป็นไพ่ใหญ่สองใบ และ rundown คือไพ่ทั้งสี่อยู่ในช่วงห้าอันดับ}
\small
\input{generated/table_archetype_rules_th.tex}
\end{table}

\begin{table}[H]
\centering
\caption{Archetype: สัดส่วนมือ, tier ของ Hwang ที่อยู่ในกลุ่ม (P, S, M, T เป็น \%) และการตัดสินใจแรกที่แนะนำแยกตาม
ตำแหน่ง}
\small
\input{generated/table_archetypes_th.tex}
\end{table}

\begin{figure}[H]
\centering
\includegraphics[width=\linewidth]{generated/fig_archetypes_th.pdf}
\caption{สัดส่วนการตัดสินใจแรกจาก solver ของแต่ละ archetype แยกตามตำแหน่ง}
\end{figure}

archetype เก้าจากสิบเอ็ดกลุ่มมีมือจาก tier ของ Hwang อย่างน้อยสอง tier ที่มีสัดส่วนตั้งแต่ 10\% ขึ้นไป กลุ่ม
``ไพ่ใหญ่ดอก nut'' เป็น Speculative ครึ่งหนึ่งและ Marginal หนึ่งในสาม กลุ่ม ``คู่ใหญ่แข็ง'' เป็น Marginal เป็นส่วนใหญ่
และ rundown ต่ำแบบ single-suited ที่ Hwang เรียกว่า Premium ตกไปอยู่ในกลุ่มมืออื่นทั้งหมด

\paragraph{การประเมิน} ตารางสุดท้ายเทียบการจัดกลุ่มสองวิธี วิธีแรกคือ EV ที่แต่ละกฎเสียเทียบกับทางที่ดีที่สุดของแต่ละมือ
โดยเลือก action ของแต่ละกลุ่มจากมือครึ่งหนึ่งแล้ววัดผลกับอีกครึ่ง วิธีที่สองคือสัดส่วนความแปรปรวนของอัตรา fold ของ
solver บนมือทั้งหมดที่การจัดกลุ่มอธิบายได้ archetype เสีย EV ราวหนึ่งในสามของคำแนะนำของ Hwang ตั้งแต่ UTG ถึง
cutoff และราวครึ่งหนึ่งที่ button และห่างจากกลยุทธ์ 780 กลุ่มของ solver เองไม่กี่ส่วนร้อยของ big blind ทุกการจัดกลุ่ม
ทำได้ไม่ดีที่ small blind ซึ่งมือส่วนใหญ่ใกล้เสมอตัว

\begin{table}[H]
\centering
\caption{การจัดกลุ่มแต่ละแบบเลียนแบบการตัดสินใจแรกของ solver ได้ดีแค่ไหน}
\small
\resizebox{\linewidth}{!}{\input{generated/table_evaluation_th.tex}}
\end{table}

\section{ข้อจำกัด}
\begin{itemize}
\item ใช้ความลึกเดียว (20 big blinds) bet ขนาด pot อย่างเดียว ไม่มี rake และ ante archetype ถูกสร้างสำหรับเกมนี้
และควรหาใหม่เมื่อ stack ลึกขึ้น ซึ่ง implied odds จะเข้าข้างมือสาย straight ของ Hwang
\item การเล่นหลัง flop ใช้การจัดกลุ่มแบบลืมประวัติ และ CFR หลายผู้เล่นไม่รับประกัน equilibrium หลักฐานด้าน
exploitability มีเพียงกำไรเล็กน้อยจากการเปลี่ยนการตัดสินใจแรก
\item การตรวจทีละมือให้ทุกคนเล่นตามกลยุทธ์ที่แก้ได้หลังการตัดสินใจแรก จึงวัดการเปลี่ยนแค่หนึ่งขั้น ไม่ใช่ best response เต็มรูปแบบ
\item archetype ถูกออกแบบหลังจากดูผล solve เดียวกับที่ใช้ประเมิน การแยกมือทดสอบป้องกันการเลือก action ของแต่ละ
กลุ่มได้ แต่ไม่ป้องกันการเลือกกฎ การทดสอบที่แข็งแรงกว่าคือ solve อิสระที่ความลึกอื่น
\item tier ที่ไม่ใช่ Premium ตรวจเพียง 400 มือต่อ tier และหลายมืออยู่ในกลุ่มก้ำกึ่งเมื่อใช้ 2{,}048 ดีลต่อมือ
\end{itemize}

\section{สรุป}
ที่ 20 big blinds tier ของ Hwang เรียงมือได้ถูก แต่วางผิดที่สองตระกูล คือมือสาย straight ไพ่ต่ำที่ไม่ใช่ double-suited
ซึ่งระบบของเขาเรียกว่า Premium หรือ Speculative และมือไพ่ใหญ่ที่มีไพ่ห้อย ซึ่งเขาเรียกว่า Marginal การตัดสินใจแรกของ
solver อธิบายได้หลัก ๆ ด้วยไพ่ใหญ่ ดอกของ A และระดับของคู่ archetype สิบเอ็ดกลุ่มที่สร้างจากปัจจัยเหล่านี้เลียนแบบ
การตัดสินใจแรกของ solver ได้ดีกว่าสี่ tier มาก และเกือบเท่ากลุ่มของ solver เอง โดยยังง่ายพอจะใช้ที่โต๊ะได้

\begin{thebibliography}{9}
\bibitem{hwang} J.~Hwang. \emph{Pot-Limit Omaha Poker: The Big Play Strategy}. Lyle Stuart
(Kensington), 2008, บทที่ 4.
\end{thebibliography}

\end{document}
